\documentclass[10pt,a4paper,roman]{moderncv}
\usepackage[utf8]{inputenc}
\moderncvstyle{classic}
\moderncvcolor{green}
\usepackage[top=1.4cm, bottom=1.5cm, left=2.0cm, right=2.0cm]{geometry}

% personal data
\name{ }{Marcus N\'obrega Gomes Junior, Ph.D.}
\address{450 Jane Stanford Way, Bldg. 320, Braun Hall}{Stanford, CA 94305}{USA}
\phone[mobile]{(+1) 520 390-2813}
\email{mngomes@stanford.edu}

\begin{document}

\recipient{Editor-in-Chief}
{Journal of Hydrology}
\date{\today}
\opening{Dear Editor-in-Chief,}
\closing{Sincerely,}
\enclosure[Postdoctoral Researcher]{Department of Earth System Science, Stanford University}

\makelettertitle
\footnotesize

I am pleased to submit the manuscript entitled ``Time Distribution of Heavy Rainfall in Brazil: Empirical Huff Curves from 290,164 Sub-Daily Storm Events'' for consideration for publication in \textit{Journal of Hydrology}.

Temporal rainfall distributions are essential for design-storm construction, urban drainage design, flood modelling, and hydrological risk assessment. In Brazil, however, engineering practice has commonly relied on the original Huff curves derived from 261 storms recorded in Illinois, USA, despite the strong climatic contrast between the temperate Midwestern United States and Brazil's tropical and subtropical rainfall regimes. This manuscript addresses that gap by deriving national-scale empirical Huff curves for Brazil using sub-daily observations from the ANA telemetric rainfall network.

The study analyzes 3,164 ANA stations and retains 290,164 qualifying heavy-rainfall events from 1,045 stations spanning 2010--2025 after applying record-length, missing-data, and event-selection criteria consistent with the original Huff framework. It provides empirical cumulative rainfall distributions and 7th-degree polynomial coefficients at station, biome, state, and municipality scales, and quantifies departures from the Huff (1967) Illinois reference using objective curve-distance metrics and uncertainty bounds.

The principal findings are directly relevant to the readership of \textit{Journal of Hydrology}. The first quartile (front-loaded storms) dominates at 94.4\% of accepted Brazilian stations, with especially strong dominance in the Amaz\^onia and Cerrado biomes. Brazilian median Q1 curves are systematically steeper than the Illinois reference, implying faster rainfall accumulation early in storm duration. A controlled SCS-CN design-hydrograph experiment shows that replacing the Illinois reference with the Brazilian empirical curves increases design peak discharge by a median of 8\%, and by up to 11\% in the Cerrado, indicating that direct transfer of the Illinois reference to Brazil can be unconservative for peak-flow estimation.

The manuscript also provides practical outputs for hydrological engineering and reproducible research. Biome-, state-, and municipality-level coefficient tables are made available as open supplementary data, together with an interactive web atlas for station metadata, dominant quartiles, empirical Huff curves, polynomial coefficients, and design-storm hyetographs generated from user-specified rainfall depth, duration, and time step.

I believe the manuscript is well aligned with \textit{Journal of Hydrology} because it combines national-scale rainfall observations, event-based storm temporal characterization, uncertainty analysis, hydrological-impact assessment, and open data products for practical design applications in a large tropical--subtropical country.

I confirm that this manuscript is original, has not been published previously, and is not under consideration elsewhere. I also confirm that I have no known competing financial interests or personal relationships that could have appeared to influence the work reported in this paper. The ANA telemetric rainfall data are publicly available, and the derived coefficient tables, code, and web atlas are available through the project repository.

Thank you for considering this manuscript for publication in \textit{Journal of Hydrology}. I look forward to your evaluation.

\makeletterclosing
\end{document}
